\subsection{An enclosure theorem is not a complexity breakthrough}
The finite-basis electronic-structure problem in the setting of
\citet{OGorman2022} is a worst-case QMA-complete ground-energy promise
problem. The relevant precision is an inverse-polynomial additive promise
gap under the paper's input and scaling conventions, not a universal relative
error guarantee for all molecular energies. Fermionic $N$-representability
has a corresponding QMA-completeness result \citep{Liu2007}.
These classifications concern families of encoded instances. They neither
prove that every molecule requires exponential time nor provide a probability
measure under which most molecules are easy. No separation of complexity
classes is assumed as a proved theorem here.

The block estimate reduces the \emph{verification statement} to a few scalar
inequalities once their hypotheses are supplied. It does not reduce the cost
of constructing those hypotheses. The active eigenproblem, the complementary
lower bound, the coupling norm, integral generation, and verified rounding all
contribute to the cost. An active-space dimension that is small relative to
the full determinant space need not be small as a function of the input size.
Moreover, a complement estimate may be cheap and so loose that the final
barrier interval is inconclusive.

The same issue arises for RDM and tensor representations. The dimensions of
an explicit $p$-RDM on antisymmetric $p$-tuples are
$\binom Kp\times\binom Kp$. For fixed $p$, this gives a polynomial
number of entries, but does not make exact $N$-representability easy. For
commuting SOS variables, if the number of variables is $m$, the number of
monomials of total degree at most $d$ is $\binom{m+d}{d}$: append a slack
exponent to the $m$ nonnegative exponents so their sum is exactly $d$, and
count weak compositions into $m+1$ parts. In particular, taking $m=K^4$
and $d=3$ gives
\[
 \binom{K^4+3}{3}
 =\frac{(K^4+3)(K^4+2)(K^4+1)}6=\Theta(K^{12}),
\]
not $O(K^6)$. This count is only for that explicitly chosen commuting
monomial basis; a fermionic operator construction must define its own basis.
A dimension count alone supplies neither an SDP runtime nor a convergence
rate, conditioning estimate, or accuracy guarantee.

\subsection{What an average-case statement would require}
An average-case complexity claim needs a family of probability distributions
$\mathcal D_n$ on finite encoded inputs, an algorithm, a success criterion,
an accuracy target, and a cost model. Qualitative labels such as stable,
local, or chemically typical do not specify a normalized distribution.
Empirical success on a selected set is evidence about that set, not a proof
of polynomial expected runtime.

\begin{proposition}[A sufficient tail condition, not an established chemical law]
\label{prop:average}
Let $T_n\ge0$ be the runtime of a specified algorithm under a specified
$\mathcal D_n$, where $n\ge1$. Suppose constants $C,C',c,d>0$ and
$k\ge0$ satisfy: outside an event of probability at most $C'e^{-cn}$,
$T_n\le Cn^k$, and on every input $T_n\le Ce^{dn}$. Then
\[
 \mathbb E[T_n]\le Cn^k+CC'e^{(d-c)n}.
\]
In particular these hypotheses imply a polynomial expectation if $c\ge d$.
\end{proposition}
\begin{proof}
Let $A_n$ be the exceptional event. Split the expectation as
$\mathbb E[T_n\boldsymbol1_{A_n^c}]+
\mathbb E[T_n\boldsymbol1_{A_n}]$.
The first term is at most $Cn^k\mathbb P(A_n^c)\le Cn^k$, and the
second is at most $Ce^{dn}\mathbb P(A_n)\le CC'e^{(d-c)n}$.
If $c\ge d$, the second bound is at most $CC'$.
\end{proof}

An exponentially small difficult subset is therefore not by itself a
polynomial expectation theorem: the runtime on it also matters. For example,
a runtime equal to $1$ outside an event of probability $e^{-cn}$ and to
$e^{dn}$ on that event has expectation
$1-e^{-cn}+e^{(d-c)n}$, which grows exponentially when $d>c$.
No distribution or tail estimate meeting these hypotheses has been
constructed here for surface chemistry. No percentages of easy molecules or
timelines for overcoming an exponential cost are inferred.

\subsection{A research protocol for a specified elementary step}
The results suggest a concrete sequence of tasks with explicit stopping
conditions. The aim is an informative interval, which may be too wide to
rank pathways; a failed ranking is not an algorithmic or physical prediction.

\begin{enumerate}
\item \textbf{Specify the physical question.}
For $\mathrm{Al_2O_3}/\mathrm{MoCl_5}$, record the substrate phase and
facet, termination and hydroxylation, slab size and boundary conditions,
coverage, precursor charge and spin, co-reactant, temperature, and reference
reservoirs. For CaF2 crystal deposition, resolve the precursor pair, the
co-reactant, the substrate phase and facet, and the reaction stoichiometry
before defining any elementary step.
Define which elementary steps and competing channels are to be compared.

\item \textbf{Fix a reproducible finite electronic target.}
Provide geometries, orbital basis and ordering, electron count, symmetry
sector, Hamiltonian coefficients, nuclear terms, and units. Record how each
geometry was obtained and what establishes a reactant minimum or a connected
transition state. A Hamiltonian at an unverified candidate geometry supports
only a statement about that geometry.

\item \textbf{Produce independently justified endpoints.}
Construct a normalized trial state and a bound on its expectation. Supply
$\underline a,\underline c,\overline b$ for the selected block split,
or a dual-feasible outer-relaxation certificate, with all feasibility and
rounding errors accounted for. Report the cost of the complement information
separately. Compare with full diagonalization only where the latter is
feasible, without treating the comparison itself as a scalable method.

\item \textbf{Bound every correction to the requested target.}
Use interval subtraction for electronic differences. Add justified bounds
for geometry, thermodynamic, environmental, and other corrections that are
needed for the chosen barrier definition. If a continuous path maximum is
required, supply between-sample control as in Proposition~\ref{prop:path}.
Do not infer a free-energy or nonadiabatic error bar from electronic solver
convergence alone.

\item \textbf{Set the precision from the kinetic decision.}
State the rate model, temperature, energy-unit convention, and positive
prefactor interval. Use Proposition~\ref{prop:log-budget} to decide whether
the energy precision can meet the desired rate ratio. Apply
Corollary~\ref{cor:ranking} only when its strict separation holds; otherwise
report an unresolved ordering and identify the dominant interval width.

\item \textbf{Define the network and observable separately.}
For a finite-state kinetic model, enumerate states, transitions, environmental
conditions, and all generator entries. Establish a rate bound for every entry
and use Corollary~\ref{cor:rates-generator}. A bounded state observable then
inherits Theorem~\ref{thm:ctmc}. Any relation to accumulated deposition or
film thickness requires its own model and error analysis. Validate physical
inputs against independently sourced data when available.
\end{enumerate}

\subsection{Open mathematical and computational questions}
The central open problem is a useful complementary-space lower bound whose
construction avoids enumerating the full omitted determinant space. Section
\ref{sec:worked} supplies the reduction that makes this a two-scalar problem, and
the certificate for it is kernel-checked. Both scalars are now measured for water
in STO-3G, and the measurement is what closes the block route rather than merely
obstructing it: $\tau$ is pinned at $\norm{O}_2$ by similarity invariance and
exceeds the diagonal gap by a factor near four. The single-reference Schur step
replaces that route and reaches $0.86$~kcal\,mol$^{-1}$, conditional on one
positive-definiteness hypothesis that the cheapest spectral tool does not certify.

Write the complement block as $H_C = D + O$ with $D$ diagonal and $O$ carrying the
off-diagonal coupling. If every diagonal entry satisfies $D_i \ge d$ and the
off-diagonal quadratic form obeys $|Q_O(x)| \le \tau\norm{x}^2$, then
\[
  \lambda_{\min}(H_C) \;\ge\; d - \tau .
\]
The chain is elementary: $d\norm{x}^2 \le \ip{x}{Dx}$ from $D_i \ge d$, and
$-Q_O(x) \le \tau\norm{x}^2$ from the coupling bound, so the sum is at least
$(d-\tau)\norm{x}^2$ for every $x$. This is formalized as
\path{complement_ground_bound} with the normalized form
\path{complement_ground_bound_normalized}, the decoupled case
\path{complement_bound_tight_when_decoupled} showing the gap equals $\tau$ exactly,
and \path{complement_bound_improvement} recording that tightening $\tau$ buys
precision one-for-one with raising $d$. The ingredients are textbook; what the
composition buys is that neither $d$ nor $\tau$ requires touching the
complement's dimension.

The open problem is therefore not sharpened but split: one half is solved and the
other half is closed.

The solved half is $d$. A lower bound on the minimum diagonal determinant energy
over the omitted sector is a combinatorial question, and it does not require
enumeration: a single-particle excitation provably raises the diagonal element, so
inspecting pair excitations suffices. This is formalized as \path{swap_identity} and
\path{diagForm_lower}, and measured for water, where the diagonal minimum is the
reference determinant and the remaining $440$ lie at least $0.227$~Ha higher. The
hypothesis can fail --- it does for O$_2$ described as a closed-shell singlet, where
an inappropriate reference pushes the diagonal minimum below the RHF energy --- and
detecting that failure here required the enumeration it was meant to avoid, which is
the honest limit of the result.

The closed half is $\tau$. It is not that no cheap estimate was found; it is that no
estimate can exist. For symmetric $O$ a similarity transform preserves the spectrum
and the $2$-norm equals the spectral radius, so no change of basis lowers $\tau$
below $\norm{O}_2$; sign structure cannot help because the quadratic form is
controlled through $\abs{O}$, and low rank and diagonal weighting fail on the
measured matrix as well. The measurements are in the paragraphs below.

What remains is the Schur scalar, and its difficulty has a different character: it
is not a bound that is too loose but a quantity whose construction requires the
block being avoided. Any upper estimate for the complement still cannot substitute
for the required lower bound.

\paragraph{The two scalars, measured where the truth is knowable.}
We measured $d$ and $\tau$ on water in STO-3G, whose $441$ determinants are few
enough to diagonalise exactly. The determinant-basis Hamiltonian used for this is
pyscf's own FCI kernel, applied to unit vectors, and the construction is checked
against \texttt{fci.FCI} with an abort on disagreement; a first hand-built version
was wrong by $33.8$~Ha and everything below supersedes it.

The exact values are
\[
d = -74.962707\ \text{Ha},\qquad
\Delta_{\mathrm{diag}} = 0.447404\ \text{Ha}\ (12.17\ \text{eV}),
\]
and the coupling constant is
\[
\tau_{\text{two-sided}} = \norm{O}_2 = 2.026458\ \text{Ha}\ (55.14\ \text{eV}),
\qquad
\tau_{\text{one-sided}} = -\lambda_{\min}(O) = 1.750272\ \text{Ha},
\]
where $\Delta_{\mathrm{diag}}$ is the gap from the lowest determinant to the next.
Note $d - E_{\mathrm{RHF}} = 1.4\times10^{-14}$~Ha: the diagonal minimum \emph{is}
the reference determinant's energy, because a determinant's diagonal element is by
definition the energy of that determinant. Of the two constants a signed
quadratic-form certificate would use the one-sided value, which is the sharper;
quoting only $2.026458$ overstates the deficit by a factor $1.16$.

Three conclusions follow. First, $d$ needs no enumeration, provided the reference
determinant is shown to be the diagonal minimum. Second,
$\tau/\Delta_{\mathrm{diag}} = 4.53$ (or $3.91$ one-sided): the off-diagonal
coupling exceeds the diagonal gap by a factor near four, so the certified bound
$d-\tau = -76.99$~Ha lies far \emph{below} the FCI energy $-75.0120$~Ha and the
certificate is sound but vacuous. Third, and this is what closes the route rather
than merely obstructing it, \emph{no change of basis can lower $\tau$}.

\paragraph{Why the block route is closed, not merely hard.}
A diagonal similarity $SOS^{-1}$ preserves the spectrum, and a symmetric matrix has
2-norm equal to its spectral radius, so
\[
  \inf_{S=\operatorname{diag}(s)>0}\norm{SOS^{-1}}_2 = \norm{O}_2 .
\]
A search over $20\,000$ random scalings confirms it: every attempt was worse
(minimum $59.33$~Ha) and none fell below $\norm{O}_2$. Sign information cannot help
either, because the quadratic form is controlled through $|O|$
(\path{abs_quadform_le_absO}), which discards all sign cancellation at the first
step, and $O$ has no usable gauge anyway: $9536$ entries are positive and $8452$
negative. Low rank is also excluded, $O$ having full rank $441$ with slowly decaying
singular values $2.03,1.76,1.75,1.70,\dots$, and $87$ of its eigenvalues lie below
$-\Delta_{\mathrm{diag}}$, so the negative part is not a small perturbation. The
weighted Collatz--Wielandt family (\path{quadform_le_diag_scaled}) is likewise
stated for $|O|$ and inherits the same floor.

Cheap proxies for $\tau$ are far from tight: the Gershgorin bound gives
$3.759$~Ha against a true norm of $2.026$~Ha, and the largest row of $O$ has $80$ of
$441$ entries nonzero, so row sparsity alone would not close the factor-of-$4.5$
deficit.

\paragraph{The residual route reaches the same wall from the other side.}
Temple's inequality
\[
  \lambda_{\min}\ \ge\ \langle\psi|H|\psi\rangle
   - \frac{\norm{(H-\langle H\rangle)\psi}^{2}}{\beta-\langle\psi|H|\psi\rangle}
\]
requires $\beta$ to be a strict lower bound on the second eigenvalue. Supplying the
\emph{true} $\lambda_2$, which is not available in practice, gives a bound within
$5.37$~kcal\,mol$^{-1}$ at a $200$-dimensional trial space, $231$ times tighter than
the operator-norm certificate. But the only computable bound at hand, Gershgorin,
gives $-75.983$~Ha, below the ground state, so the hypothesis fails and the bound
does not exist. The method relocates the difficulty from bounding $\norm{O}$ to
certifying a spectral gap.

\paragraph{A single-reference Schur step does better, and needs one hypothesis.}
Splitting about the reference determinant, $H = \begin{pmatrix} d & b^T\\
b & H_r\end{pmatrix}$, the Schur determinant identity makes $\lambda$ an
eigenvalue exactly when $\lambda = g(\lambda)$ with
\[
  g(\lambda) = d - b^T (H_r-\lambda I)^{-1} b ,
\]
defined for $\lambda<\lambda_{\min}(H_r)$. Evaluating at the Rayleigh upper bound
$\lambda_1\le d$ gives
\[
  g(d) = -75.01338150\ \text{Ha}
  \quad\text{against}\quad
  E_{\mathrm{FCI}} = -75.01201163\ \text{Ha},
\]
a gap of $0.860$~kcal\,mol$^{-1}$ rather than $1241$: a factor $1443$ improvement,
using one row of the Hamiltonian and one linear solve. Here
$H_r-dI$ has smallest eigenvalue $0.349986$~Ha, so the required shift is positive
definite.

\paragraph{Certifying the shift.}
The single hypothesis is $d < \lambda_{\min}(H_r)$, and it is certified here without
computing an eigenvalue. Gershgorin cannot do it --- and cannot ever do it, since it
tests diagonal dominance, a sufficient condition, while $A = H_r-dI$ has smallest
diagonal $0.447$~Ha against a largest row coupling of $1.380$~Ha. Diagonal dominance
therefore fails structurally, for $51$ of $440$ rows. Neither reordering nor
weighting repairs it, and we measured all of them: the Gershgorin bound is
$-1.0207$ under every one of four orderings, because it is a per-row quantity that
permutation merely relabels, and the best diagonal weighting gives $-0.7075$.
Block diagonal dominance fails by a wider margin, and monotonically worse as blocks
are refined: with $2$ blocks the in-block gap is $0.361$ against an off-block norm
of $0.903$, and with $44$ blocks it is $0.413$ against $1.831$, so the coupling
exceeds the gap by a factor of two to five at every scale. The determinant-basis
matrix has no local structure for these criteria to exploit.

A verified Cholesky factorisation does. For an approximate factor $L$ one has
$A = LL^\top + E$ with $\norm{E}_2 \le \norm{E}_F$, so by Weyl's inequality
\[
  \lambda_{\min}(A) \;\ge\; \lambda_{\min}(LL^\top) - \norm{E}_F
  \;=\; \sigma_{\min}(L)^2 - \norm{E}_F,
\]
and $\sigma_{\min}(L) = 1/\norm{L^{-1}}_2 \ge 1/(\sqrt n\,\norm{L^{-1}}_\infty)$,
which needs only two triangular solves and one residual. Measured on water,
$\norm{A-LL^\top}_F = 5.19\times10^{-14}$ and the rounding bound on forming
$LL^\top$ is $1.97\times10^{-11}$, giving $\sigma_{\min}(L)^2 \ge 7.43\times10^{-4}$:
strictly positive, hence $A \succ 0$. The certified bound is $471$ times looser than
the true $\lambda_{\min}(A) = 0.349986$, and that is immaterial --- positivity is all
the argument consumes.

With $\lambda_{\min}(A) \ge \delta > 0$ in hand, any approximate solve $\hat y$ with
residual $r = b - A\hat y$ brackets the scalar through $y - \hat y = A^{-1}r$:
\[
  b^\top A^{-1}b = b^\top\hat y + b^\top A^{-1}r,
  \qquad
  \abs{b^\top A^{-1}r} \le \norm{b}\norm{r}/\delta .
\]
Measured: $b^\top A^{-1}b = 5.067489\times10^{-2} \pm 1.82\times10^{-12}$, so
\[
  \lambda_1 \;\ge\; d - b^\top A^{-1}b \;=\; -75.01338150\ \text{Ha}
\]
is a certified lower bound at $0.860$~kcal\,mol$^{-1}$ from $E_{\mathrm{FCI}}$.
The bound is now rigorous rather than numerical.

\paragraph{What this certificate does not buy.}
It is rigorous and it is not cheaper. The Cholesky factorisation needs the whole
$440\times440$ block, at $O(n^3)$ work --- the same asymptotic order as exact
diagonalisation, with a smaller constant and no eigenvectors. The gain is in
\emph{quality}, not in cost: an unusable certificate at $1241$~kcal\,mol$^{-1}$
becomes a usable one at $0.86$, and the object that must be supplied is now named
precisely (the sign of $\lambda_{\min}(A)$ and one scalar $b^\top A^{-1}b$) instead
of being hidden in a norm. The exponential wall is untouched. Bounding
$b^\top A^{-1}b$ without forming $A$ is the obvious next target --- it is a
quadratic form $b^\top f(A)b$ with $f(t)=1/t$, accessible to Lanczos
tridiagonalisation and Gauss--Radau quadrature in $O(k\,\mathrm{nnz})$ --- but
Gauss--Radau needs a node placed below the spectrum, so it returns the same
requirement one level down. Whether that requirement can be met by an a posteriori
Ritz bound rather than by a factorisation is open here.

A cheap candidate for $d$ fails outright. Taking
$d \ge E_{\mathrm{nuc}} + 2\sum_{i\le N/2}\varepsilon_i$ gives $-36.74$~Ha, an error
of $+38.22$~Ha ($1040$~eV) and $85$ times the diagonal gap. Electron repulsion
cannot be discarded from a diagonal bound.

\paragraph{Where the reference-minimises-the-diagonal hypothesis fails.}
The observation $d = E_{\mathrm{RHF}}$ is not universal. It holds for water, HF and
N$_2$ in STO-3G to within $10^{-13}$~Ha, but fails for O$_2$ treated as a closed-shell
singlet: there the diagonal minimum lies $0.054$~Ha \emph{below} the RHF energy and
the reference determinant is only third in the diagonal spectrum, with the gap to
the next determinant collapsing to zero. The same molecule treated as its actual
triplet, with a restricted open-shell reference, restores $d = E_{\mathrm{ref}}$
exactly and gives a $0.229$~Ha gap.

The failure is therefore not evidence of a subtle integral property but of an
inappropriate reference: a single-determinant closed-shell description of a
strongly correlated open-shell ground state. Stated as a rule, the swap-based
argument of \path{swap_identity} certifies $d$ whenever the reference is the
lowest determinant, and the first thing to check for a new system is whether that
holds --- a check that costs $O(n_{\mathrm{occ}}n_{\mathrm{virt}})$
($20$ pairs against $441$ determinants for water) and that O$_2$ fails. Whether
that failure is detectable \emph{without} the enumeration is exactly the question
this measurement cannot answer, since detecting it here required diagonalising the
whole sector.

A second problem is to combine active-state improvement and lower-bound
improvement adaptively. Widening an active space may improve a variational
upper bound, but a useful algorithm must show how to maintain and evaluate
the lower certificate as the partition changes. Localized orbitals, symmetry,
or tensor compression can aid the search without proving a uniform rate of
convergence or a polynomial runtime.

A third problem is validated numerical evaluation. Floating-point eigenvalues,
SDP feasibility residuals, and rate exponentials must be enclosed with error
bounds tight enough to preserve the desired energy and rate inequalities.
The residual correction in Proposition~\ref{prop:dual} gives one possible
pattern, but its spectral residual and trace bound still have to be verified
for the concrete SDP.

Finally, uniform geometry and path errors, a defensible thermodynamic
correction budget, a complete physical event catalog, and a measured
observable mapping remain unresolved for the named ALD targets. The present
paper provides conditional implications, not solutions to these input
problems. Its next substantive result should be a fully specified example
where each hypothesis can be checked and where the final interval either
supports or demonstrably fails to support a stated chemical decision.
